\documentclass[12pt]{article}
\usepackage{amsmath, amssymb}
\usepackage{booktabs}
\usepackage{array}
\usepackage{geometry}
\usepackage{parskip}
\geometry{margin=1in}

\title{Pragmatics of ``Can You'' Questions in RSA}
\date{}

\begin{document}

\maketitle

\section{Overview}

We develop a Rational Speech Acts (RSA) model of indirect requests, using the canonical example of ``Can you pass the salt?'' The central puzzle is why a question literally asking about ability is interpreted as a request for action. We argue that the key lies in the social asymmetry: being revealed as unwilling --- as opposed to merely unable --- is socially costly, and both speaker and listener strategically exploit the ambiguity of indirect requests to manage this risk. We build the model in three layers: a literal listener $L_0$ who responds based on the literal meaning of utterances; a pragmatic speaker $S_1$ who chooses between a direct imperative and an indirect capacity question by trading off goal achievement against the social cost of exposing the listener's unwillingness; and a pragmatic listener $L_1$ who reverses this reasoning to infer the speaker's latent intention --- whether they are seeking action or information --- and responds accordingly.


\section{The Literal Listener L$_0$}

$L_0$ responds to utterances
based solely on their literal meaning and the listener's own
psychological state, without any pragmatic reasoning about what the
speaker intended.

\subsection{Listener's State}

The listener's state is fully described by two independent binary
variables:

\[
c \in \{0, 1\} \quad \text{(capacity)}, \qquad w \in \{0, 1\} \quad
\text{(willingness)}
\]

The state $(c, w)$ is known to the listener but uncertain to the
speaker, who maintains independent marginal priors $P_S(c)$ and
$P_S(w)$ with joint distribution factorizing as:

\[
P(c, w) = P(c) \cdot P(w)
\]

Factorizing the joint distribution reflects the natural assumption that whether someone \textit{can} do something is independent of whether they \textit{want} to.

\subsection{Utterances and Literal Semantics}

The speaker chooses between two utterances $\mathcal{U} = \{u_{can}, u_{imp}\}$. 
The literal semantics of each utterance determines which responses are truthfully available to the listener.

For $u_{imp}$ --- ``Pass me the salt'' --- the literal meaning is a directive to act. 
Compliance requires both capacity and willingness, and declining is the only alternative:

\[
\text{Available}(c, w, u_{imp}) = \begin{cases}
\{\text{action}\} & c = 1,\ w = 1 \\
\{\text{decline}\} & \text{otherwise}
\end{cases}
\]

For $u_{can}$ --- ``Can you pass the salt?'' --- the literal meaning is a question about capacity. Crucially, the available responses partition along $c$ alone. An able listener can truthfully answer \textit{yes} regardless of willingness; only a willing listener will perform the \textit{action}; and an unable listener must decline:

\[
\text{Available}(c, w, u_{can}) = \begin{cases}
\{\text{action},\ \text{yes}\} & c = 1,\ w = 1 \\
\{\text{yes}\} & c = 1,\ w = 0 \\
\{\text{decline}\} & c = 0
\end{cases}
\]

An unwilling-but-able listener ($c=1, w=0$) can truthfully respond \textit{yes} to $u_{can}$ --- they are answering the literal capacity question honestly --- even though this conceals their unwillingness.
This asymmetry is what gives $u_{can}$ its face-saving property, which we develop further in $S_1$.




\subsection{L$_0$'s Responses}

\paragraph{Response cost.}

At the $L_0$ level, the only cost considered is the physical effort of
producing the response, independent of utterance:

\[
\text{cost}(r) = \begin{cases}
c_a & r = \text{action} \\
c_v & r = \text{yes} \\
c_d & r = \text{decline}
\end{cases}
\]

where $c_a > c_v \approx c_d \geq 0$. 
The key asymmetry is that performing an action is generally more effortful than giving a verbal response. 
Saying \textit{yes} and saying \textit{no} are roughly equally easy. 
No social costs appear at this level --- those belong to higher levels of reasoning.

\paragraph{Utility.}

$L_0$ selects among truthful responses weighted by cost via softmax:

\[
P_{L_0}(r \mid c, w, u)
= \frac{\mathbf{1}[r \in \text{Available}(c,w,u)]
  \cdot e^{-\alpha\,\text{cost}(r)}}
  {\displaystyle\sum_{r'} \mathbf{1}[r' \in \text{Available}(c,w,u)]
  \cdot e^{-\alpha\,\text{cost}(r')}}
\]

where $\alpha \geq 0$ is the rationality parameter. The only
non-trivial case is $c=1, w=1$ under $u_{can}$, where cost favors
\textit{yes} over \textit{action} since $c_v < c_a$. All other cases
are deterministic: $L_0$ always acts under $u_{imp}$ when $c=1, w=1$,
always says \textit{yes} when $c=1, w=0$ under $u_{can}$, and always
declines when $c=0$.


\section{The Pragmatic Speaker S$_1$}

$S_1$ is the asker. Crucially, $S_1$ has a latent \textit{intention}
$\psi$ which determines what a successful outcome looks like. The
speaker chooses an utterance to maximize expected utility given that
intention, modeling how $L_0$ will respond.

\[
\psi \in \Psi = \{\text{action},\, \text{epistemic}\}
\]

\begin{itemize}
    \item $\psi = \text{action}$: the speaker wants the listener to
        perform the action
    \item $\psi = \text{epistemic}$: the speaker genuinely wants to
        learn whether the listener is capable
\end{itemize}

The speaker selects an utterance via softmax over expected utility:

\[
P_{S_1}(u \mid \psi) \propto e^{\alpha \cdot V(u, \psi)}
\]

\[
V(u, \psi) = \text{GoalAchievement}(u, \psi)
- \beta \cdot \text{SocialRejection}(u)
\]

where $\beta \geq 0$ weights how much the speaker cares about avoiding
social rejection relative to achieving their goal. We first define the
marginal response distributions used throughout:

\[
P_S(r \mid u) = \sum_{c,w} P_S(c) \cdot P_S(w)
\cdot P_{L_0}(r \mid c, w, u)
\]
\[
P_S(r \mid c, u) = \sum_w P_S(w) \cdot P_{L_0}(r \mid c, w, u)
\]

\subsection{Goal Achievement}

The two intentions warrant different utility functions since they
measure fundamentally different goals.

When $\psi = \text{action}$, the speaker wants the action to be
performed. Since $P_S(c=1) \cdot P_S(w=1)$ is constant across
utterances and drops out of the softmax, goal achievement reduces to
the log probability that $L_0$ produces an action when the listener is
willing and able:

\[
\text{GoalAchievement}(u, \text{action})
= \log P_{L_0}(\text{action} \mid c=1, w=1, u)
\]

Note that $P_{L_0}(\text{action} \mid c=1, w=1,
u_{imp}) > P_{L_0}(\text{action} \mid c=1, w=1, u_{can})$, since
$u_{imp}$ does not split probability mass between \textit{action} and
\textit{yes}.

When $\psi = \text{epistemic}$, the speaker wants to learn the
listener's capacity $c$. Goal achievement is the mutual information
between $c$ and the response $r$, which measures how much observing
the response reduces uncertainty about $c$:

\[
\text{GoalAchievement}(u, \text{epistemic})
= I(c\,;\,r \mid u)
= \sum_{r,c} P_S(c) \cdot P_S(r \mid c, u)
  \cdot \log \frac{P_S(r \mid c, u)}{P_S(r \mid u)}
\]

Under $u_{can}$, every response perfectly reveals $c$: \textit{action}
and \textit{yes} both imply $c=1$, while \textit{decline} implies
$c=0$. So $I(c\,;\,r \mid u_{can}) = H(c)$, the maximum possible.
Under $u_{imp}$, a decline pools $c=0$ and $w=0$ together, leaving
residual ambiguity about $c$, so $I(c\,;\,r \mid u_{imp}) < H(c)$.
This means $u_{can}$ is strictly better than $u_{imp}$ for epistemic
goals.

\subsection{Social Rejection}

The social rejection term captures how much a response reveals about the listener's \textit{unwillingness}. 
The intuition is that being caught in unwillingness is socially costly --- it damages the relationship between speaker and listener. 
A response that shifts the posterior over $w$ far from the prior is more costly than one that leaves it near the prior.

This is naturally formalized as the mutual information between $w$ and
$r$:

\[
\text{SocialRejection}(u) = I(w\,;\,r \mid u)
= \sum_r P_S(r \mid u)
  \cdot D_{KL}\bigl(P_S(w \mid r, u) \,\|\, P_S(w)\bigr)
\]

where the posterior over unwillingness given each response is:

\[
P_S(w=0 \mid r, u)
= \frac{\displaystyle\sum_c P_S(c) \cdot P_S(w=0)
  \cdot P_{L_0}(r \mid c, w=0, u)}{P_S(r \mid u)}
\]

The face-saving property of $u_{can}$ follows directly. Under
$u_{can}$, a decline reveals nothing about willingness since only
$c=0$ produces it --- the posterior $P_S(w=0 \mid \text{decline},
u_{can})$ equals the prior $P_S(w=0)$. The \textit{yes} response
carries some unwillingness signal, but it is diluted by willing
listeners who also say \textit{yes}. Under $u_{imp}$, a decline
concentrates unwillingness inference above the prior since $w=0$
contributes directly.

\subsection{Summary}

The full value function is:

\[
V(u, \psi) = \begin{cases}
\log P_{L_0}(\text{action} \mid c=1, w=1, u)
  - \beta \cdot I(w\,;\,r \mid u)
& \psi = \text{action} \\[6pt]
I(c\,;\,r \mid u) - \beta \cdot I(w\,;\,r \mid u)
& \psi = \text{epistemic}
\end{cases}
\]

The two utterances trade off as follows:

\vspace{0.5em}
\begin{center}
\begin{tabular}{lcc}
\toprule
 & Goal Achievement & Social Rejection \\
\midrule
$u_{imp}$ & Higher (action only) & Higher ($w=0$ directly revealed) \\
$u_{can}$ & Lower (split with \textit{yes}) & Lower ($w=0$ diluted) \\
\bottomrule
\end{tabular}
\end{center}
\vspace{0.5em}

This yields the following predictions:

\begin{itemize}
    \item When $P_S(w=1) \approx 1$, $I(w\,;\,r \mid u)$ is near zero
        under both utterances --- the speaker prefers $u_{imp}$ for its
        higher goal achievement
    \item When $P_S(w=0)$ is non-trivial and $\beta$ is high, the
        speaker prefers $u_{can}$ for its face-saving dilution of
        unwillingness inference
    \item For $\psi = \text{epistemic}$, $u_{can}$ dominates $u_{imp}$
        on both dimensions simultaneously --- higher $I(c\,;\,r \mid
        u)$ and lower $I(w\,;\,r \mid u)$ --- making utterance choice
        unambiguous
\end{itemize}

\section{The Pragmatic Listener L$_1$}

$L_1$ knows their own $(c, w)$, observes utterance $u$, and reasons
about the speaker's latent intention $\psi$ before deciding how to
respond. $L_1$'s response depends jointly on their own state, the
inferred intent, and what responses are truthfully available under $u$.

\subsection{Prior Over Speaker Intent}

Before hearing any utterance, $L_1$ forms a prior over $\psi$ by
reasoning about what kinds of intent are plausible given the context.
The key insight is that epistemic intent --- genuinely wanting to learn
about capacity --- only makes sense if the speaker is actually uncertain
about $c$. A speaker who obviously knows the listener can comply would
have no reason to ask ``can you?'' as a genuine information question.
So the prior probability of epistemic intent should scale with how
uncertain the speaker is about $c$.

This uncertainty is captured by the entropy of the speaker's prior over
capacity:

\[
H(c) = -P_S(c=1)\log P_S(c=1) - P_S(c=0)\log P_S(c=0)
\]

Normalizing by the maximum entropy $\log 2$ gives a prior in $[0,1]$:

\[
P_{L_1}(\psi = \text{epistemic}) = \frac{H(c)}{\log 2},
\qquad
P_{L_1}(\psi = \text{action}) = 1 - \frac{H(c)}{\log 2}
\]

The intuition is clean: when $P_S(c=1) \approx 1$, $H(c) \approx 0$
and epistemic intent gets prior probability near zero --- asking about
capacity you are already certain about is implausible as a genuine
question. When $P_S(c=1) = 0.5$, $H(c) = \log 2$ and the prior fully
favors epistemic intent --- maximum capacity uncertainty makes the
question entirely natural.

\subsection{Intent Inference}

$L_1$ updates the prior over $\psi$ upon observing $u$, using $S_1$ as
the generative model:

\[
P_{L_1}(\psi \mid u)
= \frac{P_{S_1}(u \mid \psi) \cdot P_{L_1}(\psi)}
  {\displaystyle\sum_{\psi'} P_{S_1}(u \mid \psi')
  \cdot P_{L_1}(\psi')}
\]

The likelihood $P_{S_1}(u \mid \psi)$ is exactly the utterance
probability computed in $S_1$. The utterance is thus informative about
intent in two complementary ways: through the likelihood, which
captures how rationally each utterance serves each intention, and
through the prior, which captures how plausible each intention is given
the context.

\subsection{L$_1$'s Response}

$L_1$ selects a response by maximizing expected goal satisfaction net
of physical cost, where the expectation is taken over the inferred
intent distribution. We build up the utility function in four steps.

\paragraph{Step 1: Goal satisfaction of Responses.}

$\text{GoalSat}(r, \psi)$ captures how well response $r$ satisfies
intent $\psi$:

\[
\text{GoalSat}(r, \psi) = \begin{cases}
1 & r = \text{action},\ \psi = \text{action} \\
1 & r = \text{action},\ \psi = \text{epistemic} \\
1 & r = \text{yes},\   \psi = \text{epistemic} \\
0 & r = \text{yes},\   \psi = \text{action} \\
1 & r = \text{decline}
\end{cases}
\]

The key structural observation is that \textit{action} satisfies both
intents simultaneously --- performing the action also demonstrates
capacity --- while \textit{yes} only satisfies $\psi =
\text{epistemic}$ and fails to deliver if the speaker actually wanted
the action. 

\paragraph{Step 2: Expected goal satisfaction.}

From the Bayesian inference in the previous subsection, $L_1$ holds a
posterior over the speaker's intent:

\[
P_{L_1}(\psi \mid u) = \frac{P_{S_1}(u \mid \psi) \cdot P_{L_1}(\psi)}
{\sum_{\psi'} P_{S_1}(u \mid \psi') \cdot P_{L_1}(\psi')}
\]

These posteriors determine how much weight each intent receives when
evaluating the expected utility of each response:

\[
\mathbb{E}_\psi[\text{GoalSat}(\text{action})]
= \sum_\psi P_{L_1}(\psi \mid u) \cdot 1 = 1
\]

This implies that \textit{action} is a dominant response at zero cost, so a costless listener should always act when able and willing.

\[
\mathbb{E}_\psi[\text{GoalSat}(\text{yes})]
= P_{L_1}(\psi = \text{epistemic} \mid u)
\]

\paragraph{Step 3: Physical cost.}

The utility of each response subtracts its physical cost:

\[
U(r \mid c, w, u)
= \log \mathbb{E}_\psi[\text{GoalSat}(r, \psi)] - \text{cost}(r)
\]

Expanding for the non-trivial case $c=1, w=1$ under $u_{can}$:

\[
U(\text{action} \mid c=1, w=1, u_{can}) = 0 - c_a = -c_a
\]
\[
U(\text{yes} \mid c=1, w=1, u_{can})
= \log P_{L_1}(\psi = \text{epistemic} \mid u_{can}) - c_v
\]

The crossover condition --- where \textit{action} and \textit{yes}
yield equal utility --- is:

\[
P_{L_1}(\psi = \text{epistemic} \mid u_{can}) = e^{c_v - c_a}
\]

So \textit{action} is preferred when $P_{L_1}(\psi = \text{epistemic}
\mid u_{can}) < e^{c_v - c_a}$, and \textit{yes} is preferred
otherwise. The threshold $e^{c_v - c_a}$ decreases as the cost gap
$c_a - c_v$ increases: a larger cost of acting raises the bar for how
confident the listener needs to be in action intent before paying it.

\paragraph{Step 4: Softmax response distribution.}

$L_1$ selects among truthfully available responses via softmax over
utility:

\[
P_{L_1}(r \mid c, w, u)
\propto \mathbf{1}[r \in \text{Available}(c, w, u)]
\cdot e^{\alpha \cdot U(r \mid c, w, u)}
\]

\paragraph{Non-compliant states.}

For non-compliant listeners the utility structure is trivial since only
one response is truthfully available. When $c=1, w=0$ under $u_{can}$,
only \textit{yes} is available --- the listener exploits the literal
semantics of $u_{can}$ as face-saving cover regardless of $\psi$.
When $c=0$, only \textit{decline} is available under either utterance.

\paragraph{Summary of key predictions.}

The response of a willing and able listener ($c=1, w=1$) under
$u_{can}$ is jointly governed by the intent posterior, the physical
cost gap, and the rationality parameter:

\begin{itemize}
    \item When $c_a - c_v$ is small, the dominant \textit{action} wins
        even under high epistemic posterior --- acting is cheap
        insurance against misreading intent
    \item When $c_a - c_v$ is large, \textit{yes} is preferred unless
        $P_{L_1}(\psi = \text{action} \mid u_{can})$ is large enough
        to clear the threshold --- strong evidence of action intent is
        required before paying the physical cost
    \item The stochasticity of the response reflects genuine uncertainty
        about intent, with the softmax producing a graded mixture of
        \textit{action} and \textit{yes} that tracks both the intent
        posterior and the cost gap
\end{itemize}

\subsection{Predictions from RSA Agents}

The model produces a clean interaction between $P_S(c)$, $(c, w)$,
and utterance:

\vspace{0.5em}
\begin{center}
\begin{tabular}{cccc}
\toprule
$(c, w)$
& $u_{can}$, high $P_S(c)$
& $u_{can}$, low $P_S(c)$
& $u_{imp}$ \\
\midrule
$(1, 1)$ & \text{action} & \text{yes}    & \text{action}  \\
$(1, 0)$ & \text{yes}    & \text{yes}    & \text{decline} \\
$(0, \cdot)$ & \text{decline} & \text{decline} & \text{decline} \\
\bottomrule
\end{tabular}
\end{center}
\vspace{0.5em}

Four dynamics are worth highlighting. First, for a willing and able
listener, $P_S(c)$ determines whether $u_{can}$ is treated as a
request or an information question: when capacity is obvious (high
$P_S(c)$), the entropy-derived prior heavily favors action intent, and
the listener acts even though the action was never literally requested.
Second, when capacity is genuinely uncertain (low $P_S(c)$), the same
listener responds with \textit{yes}, treating the utterance as a
genuine ability question. Third, the strategic \textit{yes} from an
unwilling-but-able listener is the listener-side mirror of $S_1$'s
choice of $u_{can}$: the face-saving ambiguity of the indirect form
protects both speaker and listener. Fourth, $u_{imp}$ eliminates all
flexibility --- intent is unambiguous and capacity together with
willingness fully determine the response.

\end{document}